% !TEX root = ../Main.tex
\begin{table}[htb!]
\centering
\footnotesize
\caption{Parts of the composite selection score of \Cref{Equation:Composite}. The three groups have equal weight, so return makes up one ninth of the total and does not dominate it. The ratios are those defined in \Cref{Section:PerformanceMeasurement}.}
\label{Tables:Composite}
\begin{tabularx}{\textwidth}{@{}lX@{}}
\toprule
\textbf{Group} & \textbf{Metrics, each used as a percentile rank within the candidate pool} \\
\midrule
Money & Annualised return, Calmar ratio, Sterling ratio \\
\addlinespace
Safety & Sharpe ratio, Sortino ratio, inverse maximum drawdown, inverse downside volatility \\
\addlinespace
Behaviour & Two-sidedness (the weaker exposure side as a fraction of half), regime edge over the model's own rotation null, mean holding period, longest holding period \\
\bottomrule
\end{tabularx}
\end{table}